\chapter*{Kurzfassung}
\addcontentsline{toc}{chapter}{Kurzfassung}
Parametrisierte strukturmechanische Berechnungen können bei wiederholter Auswertung durch datengetriebene Ersatzmodelle beschleunigt werden. Diese Arbeit vergleicht ein direktes mehrschichtiges neuronales Netz (MLP), eine konkrete radiale Basisfunktionsvariante (RBF), ein POD--MLP und ein POD--GPR. Ausgangspunkt sind sechs Geometrie-, Material- und Lastparameter; vorhergesagt wird das vollständige statische Verschiebungsfeld an 525 Punkten mit drei Komponenten. Der Trainingspool umfasst 768 FEM-Fälle. Für die Modellwahl werden 128 neue Validierungsfälle und für die unabhängige Bewertung 256 Interpolations- sowie 128 Extrapolationsfälle verwendet.

Bei 768 Trainingsfällen beträgt der mittlere gewichtete relative Interpolationsfehler 4,64\,\% für das MLP, 7,41\,\% für die RBF-Variante, 4,78\,\% für POD--MLP und 0,55\,\% für POD--GPR. Für die definierte Extrapolation ergeben sich 16,24\,\%, 32,67\,\%, 13,74\,\% beziehungsweise 3,85\,\%. Die vollständige Fit-Zeit des ausgewerteten POD--GPR-Hauptmodells beträgt 353,57\,s und übersteigt die neuronalen Fit-Zeiten deutlich. Ein zusätzlicher Benchmark vergleicht vollständige Surrogatvorhersagen mit FEM-Lösung und Feldinterpolation auf zwölf vorab gewählten Fällen. Die daraus abgeleiteten Amortisationswerte sind bedingte Szenarien. Große relative Extrapolationsfehler treten insbesondere bei kleinen Referenznormen auf; zudem sind absolute Abweichungen und falsche Vorzeichen einzelner neuronaler Spitzenvorhersagen nachweisbar. Die Befunde gelten für das untersuchte lineare Balkenproblem und die begrenzten Such- und Testbereiche.
